% GENERATED FILE. Edit canonical lessons, then run npm run generate:books.
% canonical-source-hash: fnv1a64:c7230cf367f03f6b
% canonical-lessons: JA-C313-kiroku, JA-C313-sakaeru, JA-C313-horobiru, JA-C313-senjitsu, JA-C313-iseki

\chapter{What Remains}
\label{ch:ja-what-remains}
\hlchaptermodality{313}

\begin{chapteropening}
\textbf{By the end of this chapter:} \emph{I can say what has come down to us from the past.}

Complete the last lesson of chapter 313: I can say what has come down to us from the past.
\end{chapteropening}

\section[\texorpdfstring{\ja{きろく} (kiroku)}{きろく (kiroku)}]{\ja{きろく} (kiroku) — a record; an account}

\label{lesson:JA-C313-kiroku}



\begin{quote}
Before the new one: say the Japanese for formerly, then the Japanese for at that time.
\end{quote}

\subsection*{You'll want to know: \ja{きろく}}
\textbf{\ja{きろく}} — \emph{kiroku} — \textquotedblleft{}a record; an account\textquotedblright{}.

\begin{quote}
\textbf{\ja{むかしの} \ja{きろくが} \ja{のこっています。}} — \emph{mukashi no kiroku ga nokotte imasu} — \textquotedblleft{}Accounts from long ago survive.\textquotedblright{}
\end{quote}

\begin{cousinweb}
\textbf{\ja{きろく}} (\emph{kiroku}) is Sino-Japanese (kango): two Chinese characters read with their on'yomi, readings borrowed from Chinese. They are \emph{ki}, \textquotedblleft{}account\textquotedblright{} (Middle Chinese \emph{kiH}) and \emph{roku}, \textquotedblleft{}record\textquotedblright{} (\emph{ljowk}) (KANJIDIC2; Middle Chinese in Baxter's transcription).
\end{cousinweb}

\begin{grammarlens}[title={What you've built}]
One more word for telling what happened.
\end{grammarlens}

\subsection*{Your turn}
\begin{itemize}
\raggedright
  \item \emph{Say it:} \emph{kiroku}
  \item \emph{Say it:} \emph{kiroku}, once more
  \item \emph{Say it:} \emph{mukashi no kiroku ga nokotte imasu} — Accounts from long ago survive.
  \item \emph{Recall:} say the Japanese for formerly, then the Japanese for at that time, then say \emph{kiroku} again
\end{itemize}

\begin{tcolorbox}[breakable,colback=teal!4,colframe=teal!35!black,title={Before you move on}]
What does \ja{きろく} mean? (\textquotedblleft{}A record; an account\textquotedblright{}.) Say it once more.
\end{tcolorbox}
\section[\texorpdfstring{\ja{さかえる} (sakaeru)}{さかえる (sakaeru)}]{\ja{さかえる} (sakaeru) — to prosper; to flourish}

\label{lesson:JA-C313-sakaeru}



\begin{quote}
Before the new one: say the Japanese for at that time, then the Japanese for a record.
\end{quote}

\subsection*{You'll want to know: \ja{さかえる}}
\textbf{\ja{さかえる}} — \emph{sakaeru} — \textquotedblleft{}to prosper; to flourish\textquotedblright{}.

\begin{quote}
\textbf{\ja{むかし} \ja{この} \ja{まちは} \ja{さかえていました。}} — \emph{mukashi kono machi wa sakaete imashita} — \textquotedblleft{}Long ago this town was prosperous.\textquotedblright{}
\end{quote}

Polite \textbf{\ja{さかえます}}, \emph{sakaemasu}; an ichidan verb. Past (polite): \textbf{\ja{まちが} \ja{さかえました。}}, \emph{machi ga sakaemashita}.

\begin{cousinweb}
\textbf{\ja{さかえる}} (\emph{sakaeru}) is a native Japanese word (wago), inherited rather than borrowed from Chinese; its kanji was attached later as a kun'yomi (JMdict).
\end{cousinweb}

\begin{grammarlens}[title={What you've built}]
One more word for telling what happened.
\end{grammarlens}

\subsection*{Your turn}
\begin{itemize}
\raggedright
  \item \emph{Say it:} \emph{sakaeru}
  \item \emph{Say it:} \emph{sakaeru}, once more
  \item \emph{Say it:} \emph{mukashi kono machi wa sakaete imashita} — Long ago this town was prosperous.
  \item \emph{Recall:} say the Japanese for at that time, then the Japanese for a record, then say \emph{sakaeru} again
\end{itemize}

\begin{tcolorbox}[breakable,colback=teal!4,colframe=teal!35!black,title={Before you move on}]
What does \ja{さかえる} mean? (\textquotedblleft{}To prosper; to flourish\textquotedblright{}.) Say it once more.
\end{tcolorbox}
\section[\texorpdfstring{\ja{ほろびる} (horobiru)}{ほろびる (horobiru)}]{\ja{ほろびる} (horobiru) — to perish; to die out}

\label{lesson:JA-C313-horobiru}



\begin{quote}
Before the new one: say the Japanese for a record, then the Japanese for to prosper.
\end{quote}

\subsection*{You'll want to know: \ja{ほろびる}}
\textbf{\ja{ほろびる}} — \emph{horobiru} — \textquotedblleft{}to perish; to die out\textquotedblright{}.

\begin{quote}
\textbf{\ja{ふるい} \ja{くには} \ja{さいがいで} \ja{ほろびました。}} — \emph{furui kuni wa saigai de horobimashita} — \textquotedblleft{}The old kingdom perished in a disaster.\textquotedblright{}
\end{quote}

Polite \textbf{\ja{ほろびます}}, \emph{horobimasu}; an ichidan verb. Past (polite): \textbf{\ja{その} \ja{くにが} \ja{ほろびました。}}, \emph{sono kuni ga horobimashita}.

\begin{cousinweb}
\textbf{\ja{ほろびる}} (\emph{horobiru}) is a native Japanese word (wago), inherited rather than borrowed from Chinese; its kanji was attached later as a kun'yomi (JMdict).
\end{cousinweb}

\begin{grammarlens}[title={What you've built}]
One more word for telling what happened.
\end{grammarlens}

\subsection*{Your turn}
\begin{itemize}
\raggedright
  \item \emph{Say it:} \emph{horobiru}
  \item \emph{Say it:} \emph{horobiru}, once more
  \item \emph{Say it:} \emph{furui kuni wa saigai de horobimashita} — The old kingdom perished in a disaster.
  \item \emph{Recall:} say the Japanese for a record, then the Japanese for to prosper, then say \emph{horobiru} again
\end{itemize}

\begin{tcolorbox}[breakable,colback=teal!4,colframe=teal!35!black,title={Before you move on}]
What does \ja{ほろびる} mean? (\textquotedblleft{}To perish; to die out\textquotedblright{}.) Say it once more.
\end{tcolorbox}
\section[\texorpdfstring{\ja{せんじつ} (senjitsu)}{せんじつ (senjitsu)}]{\ja{せんじつ} (senjitsu) — the other day}

\label{lesson:JA-C313-senjitsu}



\begin{quote}
Before the new one: say the Japanese for to prosper, then the Japanese for to perish.
\end{quote}

\subsection*{You'll want to know: \ja{せんじつ}}
\textbf{\ja{せんじつ}} — \emph{senjitsu} — \textquotedblleft{}the other day\textquotedblright{}.

\begin{quote}
\textbf{\ja{せんじつ、ふるい} \ja{ともだちに} \ja{あいました。}} — \emph{senjitsu, furui tomodachi ni aimashita} — \textquotedblleft{}The other day I met an old friend.\textquotedblright{}
\end{quote}

\begin{cousinweb}
\textbf{\ja{せんじつ}} (\emph{senjitsu}) is Sino-Japanese (kango): two Chinese characters read with their on'yomi, readings borrowed from Chinese. They are \emph{sen}, \textquotedblleft{}previous\textquotedblright{} (Middle Chinese \emph{sen}) and \emph{jitsu}, \textquotedblleft{}day\textquotedblright{} (\emph{nyit}) (KANJIDIC2; Middle Chinese in Baxter's transcription).
\end{cousinweb}

\begin{grammarlens}[title={What you've built}]
One more word for telling what happened.
\end{grammarlens}

\subsection*{Your turn}
\begin{itemize}
\raggedright
  \item \emph{Say it:} \emph{senjitsu}
  \item \emph{Say it:} \emph{senjitsu}, once more
  \item \emph{Say it:} \emph{senjitsu, furui tomodachi ni aimashita} — The other day I met an old friend.
  \item \emph{Recall:} say the Japanese for to prosper, then the Japanese for to perish, then say \emph{senjitsu} again
\end{itemize}

\begin{tcolorbox}[breakable,colback=teal!4,colframe=teal!35!black,title={Before you move on}]
What does \ja{せんじつ} mean? (\textquotedblleft{}The other day\textquotedblright{}.) Say it once more.
\end{tcolorbox}
\section[\texorpdfstring{\ja{いせき} (iseki)}{いせき (iseki)}]{\ja{いせき} (iseki) — ruins; remains}

\label{lesson:JA-C313-iseki}



\begin{quote}
Before the new one: say the Japanese for to perish, then the Japanese for the other day.
\end{quote}

\subsection*{You'll want to know: \ja{いせき}}
\textbf{\ja{いせき}} — \emph{iseki} — \textquotedblleft{}ruins; remains\textquotedblright{}.

\begin{quote}
\textbf{\ja{やまの} \ja{なかで} \ja{ふるい} \ja{いせきが} \ja{みつかりました。}} — \emph{yama no naka de furui iseki ga mitsukarimashita} — \textquotedblleft{}Old ruins were found in the mountains.\textquotedblright{}
\end{quote}

\begin{cousinweb}
\textbf{\ja{いせき}} (\emph{iseki}) is Sino-Japanese (kango): two Chinese characters read with their on'yomi, readings borrowed from Chinese. They are \emph{i}, \textquotedblleft{}leave behind\textquotedblright{} (Middle Chinese \emph{ywij}) and \emph{seki}, \textquotedblleft{}tracks\textquotedblright{} (\emph{tsjek}) (KANJIDIC2; Middle Chinese in Baxter's transcription).
\end{cousinweb}

\begin{grammarlens}[title={What you've built}]
One more word for telling what happened.
\end{grammarlens}

\subsection*{Your turn}
\begin{itemize}
\raggedright
  \item \emph{Say it:} \emph{iseki}
  \item \emph{Say it:} \emph{iseki}, once more
  \item \emph{Say it:} \emph{yama no naka de furui iseki ga mitsukarimashita} — Old ruins were found in the mountains.
  \item \emph{Recall:} say the Japanese for to perish, then the Japanese for the other day, then say \emph{iseki} again
\end{itemize}

\begin{tcolorbox}[breakable,colback=teal!4,colframe=teal!35!black,title={Before you move on}]
What does \ja{いせき} mean? (\textquotedblleft{}Ruins; remains\textquotedblright{}.) Say it once more.
\end{tcolorbox}
